\section{引言}\label{sec:intro}

由 L\"uscher~\cite{Luscher:2010iy} 提出并经 L\"uscher 与
Weisz~\cite{Luscher:2011bx} 进一步形式化的梯度流形式，已在格点 QCD 的
诸多方面被证明十分有用。它的一个主要优点是：在有限流时间 $t$ 处的复合算符，
无需超出所涉参数与场之外的超紫外（\abbrev{UV}）重正化。这意味着各算符在
重正化群演化下互不混合，从而使来自不同正规化方案的结果可以特别简单地组合
起来。这一特性为格点计算与微扰计算的相互促进开辟了诱人的前景，例如对
$\alpha_s(M_Z)$ 的可能格点测定~\cite{Harlander:2016vzb}。

展示这种可能的相互促进的一个特别有力的途径，是考察复合算符在小流时间极限下
的展开：它把流动算符表达为 $t=0$ 处的 QCD 算符，并带有依赖于 $t$ 的
Wilson 系数~\cite{Luscher:2011bx}。Makino 与
Suzuki~\cite{Suzuki:2013gza,Makino:2014taa} 曾用该方法导出了能动量张量
（\emt）$T_{\mu\nu}$ 的一个与正规化无关的公式，且已得到富有前景的结果
（例如参见
\citeres{Asakawa:2013laa,Kitazawa:2016dsl,Taniguchi:2016ofw,Kitazawa:2017qab,
 Yanagihara:2018qqg}）。

\citere{Makino:2014taa} 中 \abbrev{EMT} 公式里出现的普适 Wilson 系数，
已在微扰论中算到次领头阶（\nlo）~\cite{Suzuki:2013gza,Makino:2014taa}。
就其只涉及对单个 $D$ 维动量的积分而言，这相当于一个单圈计算。本文将把该
计算推进到下一个微扰阶。\footnote{注意我们工作于无穷大体积极限。有限体积
  效应的纳入需要不同技术，例如数值随机微扰论，见
  \citere{DallaBrida:2017tru}。} 需要在此指出：梯度流形式中出现的积分属于
比常规 QCD 更一般的类型。它们含有额外的指数因子，这些因子依赖于圈动量与外
动量以及流时间变量，其中一些流时间变量还要被积分。尽管如此，首个双圈结果
早已在 \citere{Luscher:2010iy} 中以解析形式给出。而推广到三圈水平则离不开
计算机代数与数值工具的大力辅助\,\cite{Harlander:2016vzb}。从量子场论的角度
看，那条路径紧随 \citere{Luscher:2010iy} 的步骤，借助 Wick 定理把格林函数
直接表示成积分。这些积分本身用扇形分割
法\,\cite{Binoth:2003ak,Smirnov:2013eza}求值，以分离出 $D-4$ 极点，其系数
通过高精度数值方法确定~\cite{genz,holodborodko}。

在当前的计算中，我们采用一套完全独立的框架。一方面，它把梯度流形式表述为
一个五维量子场论\,\cite{Luscher:2011bx}，由此给出定义良好、尽管非标准的费曼
规则。另一方面，我们不再对所得积分作数值求值，而是用 Chetyrkin 与 Tkachov
的分部积分方法\,\cite{Chetyrkin:1981qh}把它们约化为若干主积分。这将使
\emt\ Wilson 系数的 \nlo\ 计算归结为单个不含流时间积分的单圈积分。
次次领头阶（\nnlo）的计算则给出四个不含流时间积分的双圈主积分和两个带单个
流时间积分的双圈主积分。对于一般的时空维数 $D$，所有主积分都能用标准方法
解析求出。

通过适当重正化，可以把 \emt\ 的 Wilson 系数定义为形式上不依赖重正化标度。
对固定阶微扰结果而言，这意味着重正化标度依赖性在形式上属于更高阶。于是人们
可以通过 Wilson 系数围绕某个特定``中心''值的剩余 $\mu$ 依赖来估计其微扰
不确定性。基于解析结果的形式，我们论证了该中心值的一个具体选取。我们的数值
研究表明，更高阶项确实显著减小了 $\mu$ 变差。然而，逐阶比较显示：按常规微扰
QCD 计算的习惯做法、在 $\mu\in[\mu_0/2,2\mu_0]$ 内变分而得的不确定性估计，
也许过于乐观。

虽然我们把 \emt\ Wilson 系数的 \nnlo\ 表达式视为主要结果，但我们的计算还
带来一些可能在更广语境中有用的附加结果。其中包括流动夸克场重正常数
$Z_\chi$，以及构成常规（非流动）QCD 能动量张量的那组算符到 \nnlo\ 的反常
量纲矩阵。

本文其余部分结构如下：第\sct{sec:pert}节简要介绍微扰梯度流形式以确定记号；
第\sct{sec:emt}节概述 \citeres{Suzuki:2013gza,Makino:2014taa} 用该形式定义
\emt\ 的做法；计算的技术细节在第\sct{sec:calc}节描述；第\ref{sec:coefs}节
给出主要结果，即 \msbar\ 方案下到 \nnlo\ QCD 的 Wilson 系数。正如
\citere{Makino:2014taa} 所指出，\emt\ 的迹反常为该计算提供了一个受欢迎的
检验；我们在第\sct{sec:traceanom}节简述由此得到的系数函数之间关系的推导。
最后，在第\sct{sec:zij}节中，我们利用流动算符的有限性条件，推导常规 QCD 中
出现在 \emt\ 内的那组算符的反常量纲矩阵。第\ref{sec:conclusions}节给出结论。
